% Overleaf: paste into main.tex and compile with pdfLaTeX.
% Self-contained: no external figures or bibliography file required.
\documentclass[11pt,a4paper]{article}
\usepackage[margin=2.3cm]{geometry}
\usepackage[T1]{fontenc}
\usepackage[utf8]{inputenc}
\usepackage{lmodern}
\usepackage{microtype}
\usepackage{amsmath,amssymb,booktabs,tabularx}
\usepackage{tikz}
\usetikzlibrary{positioning,arrows.meta}
\usepackage[colorlinks=true,linkcolor=blue!50!black,urlcolor=blue!60!black]{hyperref}
\newcommand{\file}[1]{\nolinkurl{#1}}
\setlength{\parindent}{0pt}
\setlength{\parskip}{5pt}
\setlength{\emergencystretch}{2em}
\renewcommand{\arraystretch}{1.15}
\title{From Semantic Induction Heads to Multi-hop Circuits\\
\large Research Progress and Experiment Record}
\author{Daniella Simonovsky and Max Golubkov}
\date{Working record --- updated after the OLMo-2-7B comparison (September 2026)}

\begin{document}
\maketitle

\section{Scope and Current Status}
Research log for the project, from proposal to the model-platform decision. It is a record for writing the final report, not the report itself.

\textbf{Completed:} literature study and methodology design; synthetic multi-hop datasets with built-in controls (four major revisions); evaluation infrastructure (multi-GPU Slurm, resumable runs, manifests, integrity checks); behavioral pilots on Pythia-1B and Pythia-6.9B; an order-pair analysis that separates content-based from layout-based success; a tokenizer-matched comparison run on OLMo-2-7B with a pre-registered decision rule; and the platform decision: \textbf{OLMo-2-7B is the primary model for the mechanistic phase, Pythia-6.9B is retained as a contrast model}.

\textbf{Not yet done:} Relation Index (RI) measurements on our models, activation or path patching, circuit extraction, faithfulness evaluation, checkpoint sweeps. No Semantic Induction Circuit (SIC) has been identified yet.

\section{Research Question}
From the proposal (\emph{From Heads to Circuits: Identifying Semantic Induction Circuits for Multi-hop Reasoning in LLMs}):
\begin{quote}
How do language models process multi-hop semantic relations in context? Are these relations handled by individual semantic induction heads or by distributed circuits across layers?
\end{quote}
Concretely: identify the circuit responsible for two-hop in-context composition (e.g.\ $A\!\to\!B$, $B\!\to\!C \Rightarrow A\!\to\!C$ over invented entities), validate it causally, and track its emergence across public training checkpoints. The behavioral phase below established \emph{which model} reliably performs the task and \emph{what} the circuit must explain (including its characteristic failures); the edge-rewiring conditions double as the corrupted inputs that patching requires.

\section{Literature and Method Stack}
Core sources: the transformer-circuits framework (residual stream, QK/OV, head composition) \cite{framework}; induction heads and their role in ICL \cite{olsson}; semantic induction heads and the Relation Index, the direct starting point \cite{ren}. Causal methodology: the IOI circuit and path patching \cite{ioi}; activation-patching practice \cite{patchguide}. For circuit discovery at scale we selected position-aware edge attribution (PEAP-style schemas suit our templated prompts) with gradient-based edge scoring (EAP-IG) for screening, manual activation/path patching for verification, and a pre-registered faithfulness criterion. Two standing cautions from the reading: an SIH label does not establish the classical copying circuit (separate prefix-matching/copying controls are planned), and RI is observational --- a screening signal, not causal evidence.

\section{Planned Mechanistic Pipeline}
\begin{enumerate}
\item \textbf{Behavioral validation} (done): direct retrieval, two-hop composition, edge rewiring, robustness; order-pair design.
\item \textbf{Candidate screening}: RI (with random-target null), attention patterns, per-component logit attribution. Compute order matters: apply the QK dominance filter from cached attention first, and project through OV to the vocabulary only for surviving (head, position) pairs at entity positions --- the naive all-heads $\times$ all-positions OV projection costs $\sim$30 forward passes per prompt.
\item \textbf{Node interventions}: activation patching between clean and rewired prompts; effect measured on the answer margin
\[
m(x)=\log P(y_{\mathrm{target}}\mid x)-\log P(y_{\mathrm{alternative}}\mid x),
\]
with the three-way logit contrast (correct new answer / new bridge / old answer) as the finer-grained outcome.
\item \textbf{Directed connections}: path patching to substantiate ``$h_1$ feeds $h_2$''; Q/K vs V effects separated.
\item \textbf{Circuit validation}: faithfulness of the sparse subgraph (recovery threshold fixed in advance), held-out families, generalization across the two fact orders.
\item \textbf{Checkpoint sweep}: repeat screening plus a fixed patching battery on $\sim$10--15 OLMo-2-7B revisions; a \emph{costing dry-run} on the final checkpoint (time per cached forward, per full RI pass, per patched forward; backward-pass memory fit for gradient attribution on 11--12\,GB GPUs) precedes any sweep commitment.
\end{enumerate}

\section{Models, Data, and Infrastructure}
\textbf{Models} (base LMs, no fine-tuning; ICL only; fp16; revisions pinned in per-model lock files):
\begin{itemize}
\item Pythia-1B: \file{f73d7dcc545c8bd326d8559c8ef84ffe92fea6b2};
\item Pythia-6.9B: \file{c0e3eee36dc47af0c49f361c74cfe459c09f7f23};
\item OLMo-2-1124-7B: \file{7df9a82518afdecae4e8c026b27adccc8c1f0032} (weights stored fp32, 27.2\,GiB on disk; run in fp16, $\approx$14.6\,GiB across two 11--12\,GB GPUs --- the same infrastructure as Pythia-6.9B).
\end{itemize}
Compute: TAU Slurm (\texttt{studentkillable}); node \texttt{s-002} excluded (faulty GPU~0).

\textbf{Data design.} Synthetic families with invented names: a query chain and a distractor chain per family (kinship anchor world: child $\to$ mother $\to$ grandmother; breadth worlds: containment, location). Gold relations known by construction; an independent verifier re-solves every question from prompt text alone. Controls built into the data: mention balance (frequency rule scores exactly 50\%), candidate token-length matching, no options line (candidates scored outside the prompt), demo/test name disjointness, and two fact-order variants per question.

\textbf{Runner.} \file{pilot_v2/pilot.py} + \file{parallel.py}: multi-GPU sharded evaluation by family, resumable, manifests with model/data/code hashes. Now multi-model (\texttt{--model \{pythia-6.9b, olmo2-7b\}}); each dataset row records the tokenizer(s) it was length-matched for, and the runner \emph{refuses} a model/dataset combination that was not audited together. Numerical setup: fp16 weights, math-only SDPA, explicit finite-value checks (adopted after non-finite logits in early 6.9B runs).

\section{Evaluation Measures}
\textbf{Candidate accuracy}: the correct name outscores one matched distractor by summed log-probability of the full continuation (leading space and period included; tokenization boundary asserted). \textbf{Free accuracy}: the first generated name equals gold. \textbf{Order-pair (both-orders) accuracy} --- the primary robustness measure: each question appears in two versions differing only in whether the query family's facts are far from or adjacent to the question; a question counts only if answered correctly in \emph{both}. Chance is 25\%, and a model that answers by position fails at least one version by construction. Families are the statistical unit for interval estimates. Free-generation errors are categorized (bridge entity / distractor / other) rather than only counted.

\section{Experiment History (Pythia)}
\subsection{Yes/No formats: a dead end (Pythia-1B, then 6.9B)}
The initial entailment task (Yes/No over graph-derived facts) failed at both scales: the 1B model was label-biased (88\% on positives, 19\% on negatives); the 6.9B model answered ``No'' almost universally (1{,}154--1{,}200 of 1{,}200) across two demonstration revisions, including a balanced one (\texttt{prompt-v3}). Conclusion, consistent with \cite{ren} (App.~D): base LMs of this scale do not comply with verification formats. All later work uses name completion.

\subsection{Completion-v1: real one-hop signal, confounded two-hop}
Candidates shown in the prompt; 1{,}200 questions per shot condition. Direct retrieval reached 91\% (candidate, 4-shot), but the two-hop result was uninterpretable: the correct candidate was usually the more frequently mentioned one (a frequency rule scored 100\% on basic two-hop), candidate position in the options line biased answers, and half of the free two-hop answers named the bridge entity. Changed-fact pair success was 3--19\%. Verdict: test flaws, not measured ability.

\subsection{Completion-v2: over-corrected into a negative result}
Rephrasing as explicit graph navigation (``shortest chain of exactly $k$ steps'') with fully balanced candidates collapsed \emph{everything} to $\approx$50\%, including one-hop. Retained as a documented negative result: base LMs are highly sensitive to unnatural, meta-linguistic prompt formats. Fix confounds in data structure and scoring; keep surface language natural.

\subsection{Completion-v3 / v3.1: the current instrument}
Natural language, no preamble, no options line; the composed relation has its own name (\emph{grandmother}), so the bridge is never a valid answer in the anchor world; mention balance and token-length matching by construction; 100 families (68 kinship, 16 containment, 16 location) $\times$ 6 questions $\times$ 2 fact orders = 1{,}200 prompts per shot condition (0/4/12).

v3's first run contained a demonstration confound we caught by the order split: demos coupled hop count to fact-block order and taught a position rule (two-hop 34--50\% in one order vs 93--97\% in the other). \textbf{v3.1} fully crosses hops $\times$ order in the demonstration bank. All headline results below are v3.1.

\textbf{Pythia-6.9B on v3.1 --- two regimes} (kinship, candidate accuracy, order$_{\mathrm{far}}$/order$_{\mathrm{adjacent}}$/both, 12-shot):
direct $0.84/1.00/0.84$; two-hop $0.72/1.00/0.72$ (above chance in both orders, but fully asymmetric); \emph{reorder} (shuffled facts) $0.51/0.41/0.25$ --- exactly chance; \emph{broken/first} (first edge rewired to the other family's mother) $0.18/0.00/0.00$ --- systematically wrong, dominated by the \emph{old} grandmother; \emph{broken/second} $0.68/1.00/0.68$ --- tracked correctly. Interpretation (hypothesis): hop-2 (bridge$\to$answer) is entity-based; hop-1 (query$\to$bridge) is resolved by layout/locality, which the canonical layout happens to align with truth.

\section{Model Comparison: OLMo-2-7B}
\subsection{Selection}
To test whether the layout-bound regime is a training-data limitation (Pythia: $\sim$0.3T tokens) rather than a scale limitation, we compared at matched size against OLMo-2-1124-7B ($\sim$4T tokens; same 32 layers $\times$ 32 heads). Larger OLMo variants were rejected: 13B ($\approx$27\,GB fp16) needs 3--4 GPUs per job and 32B ($\approx$64\,GB) is infeasible on the partition; TransformerLens supports the 1B and 7B OLMo-2 models but not 13B/32B; and the analysis phase consists of many small jobs, where per-job GPU count dominates queue time. OLMo-2 also publishes intermediate checkpoints (\texttt{stepXXX-tokensYYYB} revisions), which the later sweep requires.

\subsection{Tokenizer audit and adapted dataset}
v3.1 candidate pairs were length-matched under Pythia's tokenizer; under OLMo-2's ($\sim$100K vocabulary) 1{,}450 of 3{,}640 rows (62 families) broke --- a direct scoring bias if uncorrected. \file{pilot_v3/adapt_data_olmo2.py} renamed 81 entities in 60 families (anchored on each family's two-hop gold; replacements length-matched under \emph{both} tokenizers and absent from all prompts). Everything else --- ids, question structure, demonstration blocks (byte-identical), mention balance --- is unchanged. Verbatim controls use demonstration names and are flagged as unmatched by design. Verified: pair-length and boundary audits clean under both tokenizers; logic verifier 0 errors on 3{,}600 questions and 19{,}200 demonstration blocks. A preflight script re-ran the audit on the cluster with the downloaded snapshot tokenizer before any GPU job.

\subsection{Results}
Pipeline calibration (OLMo controls): verbatim 20/20 candidate; stated 15/20, all five failures selecting the more \emph{recent} distractor fact over the explicitly stated answer (model behavior --- recency competition --- not a pipeline defect). A Pythia run of the same controls file is still pending.

Kinship, candidate accuracy, order$_{\mathrm{far}}$/order$_{\mathrm{adjacent}}$/\textbf{both}, 12-shot:
\begin{center}\small
\begin{tabular}{lcc}
\toprule
Check & Pythia-6.9B & OLMo-2-7B\\\midrule
direct & $0.84/1.00/\mathbf{0.84}$ & $1.00/0.90/\mathbf{0.90}$\\
two-hop & $0.72/1.00/\mathbf{0.72}$ & $0.82/0.88/\mathbf{0.78}$\\
reorder & $0.51/0.41/\mathbf{0.25}$ & $0.62/0.57/\mathbf{0.37}$\\
broken/first & $0.18/0.00/\mathbf{0.00}$ & $0.15/0.44/\mathbf{0.09}$\\
broken/second & $0.68/1.00/\mathbf{0.68}$ & $0.94/0.91/\mathbf{0.85}$\\
\bottomrule
\end{tabular}
\end{center}
Three findings:
\begin{enumerate}
\item \textbf{Order symmetry restored.} OLMo's two-hop is $0.82/0.88$ vs Pythia's $0.72/1.00$: success no longer depends on where the query family sits.
\item \textbf{Above chance on shuffled facts --- first such evidence.} Reorder both-orders: OLMo $0.41$ [Wilson $0.30,0.53$] at 4-shot ($p=2.5\times10^{-3}$ vs chance $0.25$) and $0.37$ [$0.26,0.49$] at 12-shot; Pythia is exactly at chance ($p=.75/.55$).
\item \textbf{Same broken/first failure rate, different failure.} Free-generation anatomy (kinship, 68 questions): Pythia's dominant error is the \emph{old} grandmother (rewired edge ignored --- the layout answer); OLMo's is the \emph{new} bridge (rewired edge followed by content, composition incomplete). Shot dependence for OLMo: correct-new-answer 33 (0-shot) $\to$ 13 (4) $\to$ 6 (12); bridge 25 $\to$ 38 $\to$ 41. \textbf{Well-formed demonstrations actively suppress OLMo's completion of the rewired chain} --- a demo-induced prior competing with in-context evidence.
\end{enumerate}
Caveats: broken cells are $n{=}34$ order pairs; per-world two-hop (both-orders, 12-shot) is kinship P$.72$/O$.78$, location P$.12$/O$.69$, but containment P$.44$/O$.25$ --- OLMo is \emph{worse} there (format interaction; containment must not be pooled with kinship); OLMo is somewhat more distractible in the adjacent order (distractors variant $0.72$ vs Pythia $0.91$).

\subsection{Pre-registered decision rule and outcome}
Fixed before any OLMo number: adopt OLMo iff (1) direct and two-hop at least Pythia's level; (2) broken/first above 50\% in both orders; (3) broken/second at least Pythia's level; secondary: reorder above chance. Outcome at 12-shot: (1) pass, (2) \textbf{fail}, (3) pass, secondary pass --- a partial pass. \textbf{Documented deviation:} the rule's failure clause assumed a Pythia-like (layout) failure; the observed failure is a composition bottleneck on top of working content-based binding --- the mechanism the project needs. Decision: OLMo-2-7B primary, Pythia-6.9B retained as a mechanistic contrast (two regimes on identical data strengthen, not weaken, the comparison).

\subsection{Candidate mechanisms for the demonstration-suppression effect}
Status: fact = the output distribution over \{new answer, new bridge, old answer\} shifts causally with shot count; everything below is hypothesis. (i) \emph{Truncated composition}: hop-1 succeeds on the rewired edge; the completing fact (far in one order) is not retrieved --- supported by bridge answers concentrating where that fact is far (25 vs 16 of 34). (ii) \emph{Routine confusion}: demos contain both one-hop and two-hop questions; the bridge is the correct answer to the wrong (one-hop) question. Cheap discriminating experiment: demo banks of only-mother vs only-grandmother questions on the broken/first subset. (iii) \emph{Linked-name copying}: copy the name directly linked to the queried entity (present already at 0-shot, amplified by demos); mechanistically natural for induction-type heads, testable by logit attribution. For the checkpoint sweep, the pre-registered discriminating measure is
\[
\Delta(\text{ckpt}) \;=\; \mathrm{Acc}^{\text{broken/first}}_{\text{0-shot}}-\mathrm{Acc}^{\text{broken/first}}_{\text{12-shot}},
\]
compared against the demo-benefit curve (schema mechanisms predict co-emergence with ICL use) and the two-hop curve (the ``immature ability'' account predicts shrinkage as two-hop stabilizes).

\section{Known Limitations and Audit Status}
\begin{itemize}
\item \texttt{paired\_base\_id} references omit the order suffix and do not match result ids; the shipped paired-metric files are therefore empty. Order-pair analysis was computed independently via \texttt{option\_pair\_id} (used for all both-orders numbers above); the field should be repaired in the next data revision.
\item In 4-/12-shot runs the two order variants receive different sampled demonstrations; the both-orders metric tolerates this, but single-order contrasts do not isolate fact order alone.
\item v3.1 lacks a direct+reorder control (is \emph{one}-hop layout-robust?); planned for the next revision, alongside a containment-format fix.
\item Broken-cell sample sizes ($n{=}34$ pairs) limit precision; the distractor-placement generator uses Python's \texttt{hash()}, so the exact generated files (preserved, hashed in manifests) are the reproducibility anchor.
\item Cross-model comparison relies on the dual-tokenizer-matched dataset; legacy Pythia numbers come from the original v3.1 files (40 of 100 families textually identical; the 60 renamed families preserve structure, ids, and balance by construction).
\end{itemize}

\section{Lessons and Next Steps}
Lessons: validate the task before interpreting the model (three generations of confounds were caught by built-in baselines, order pairs, and an independent logic verifier); error \emph{anatomy} distinguishes mechanisms where accuracy alone cannot (identical broken/first accuracy, opposite failure modes); pre-registered rules work, including honestly reporting a partial pass and a justified deviation; tokenizer-dependent scoring must be enforced in code, not by convention.

Next: (1) costing dry-run on the OLMo final checkpoint (doubles as the patching tooling shakedown); (2) component screening --- RI with nulls + logit attribution on the three-way contrast; (3) activation/path patching on the three pre-registered contrasts: clean vs broken/first, canonical vs reorder, 0-shot vs 12-shot broken/first; (4) the demo-composition behavioral experiment; (5) circuit extraction + faithfulness on held-out families; (6) checkpoint sweep with $\Delta$, demo-benefit, and two-hop curves; (7) run the Pythia controls file; repair \texttt{paired\_base\_id}.

\appendix
\section{Source Files}
Paths relative to \texttt{final proj/project}; raw prompts and results are preserved alongside summaries. Repo: \url{https://github.com/Maxgolu/NLP_proj}.
\begin{itemize}
\item Proposal: \file{../NLP Projects_Proposal.pdf}; papers: \file{../papers/}; plan: \file{data_pipeline_plan.tex}.
\item Runner (multi-model) and locks: \file{pilot_v2/pilot.py}, \file{parallel.py}, \file{model_lock.json}, \file{model_lock_olmo2.json}; OLMo operations and decision rule: \file{pilot_v2/README_OLMO2.md}; preflight: \file{pilot_v2/preflight_olmo2.py}.
\item Data: generator \file{pilot_v3/build_data_v3.py}; OLMo adaptation + audits: \file{pilot_v3/adapt_data_olmo2.py}, \file{pilot_v3/audit_tokenizer.py}, \file{pilot_v3/verify_logic.py}, \file{pilot_v3/data_olmo2/} (incl.\ \file{adapt_audit_olmo2.json}).
\item Yes/No era: \file{results/full_4shot_review/}, \file{results/full_sdpa_review/}, \file{results/prompt_v3_review/}; completion-v1: \file{results/completion_review/}, audit \file{results/completion_analysis/}; v2 (negative result): \file{results/completion_v2_*}.
\item v3/v3.1 (Pythia): \file{results/completion_v3_4shot_v1/}, \file{results/completion_v3_12shot_v1/}, \file{results/completion_v3_1_0shot/}, \file{results/completion_v3_1_4shot/}, \file{results/completion_v3_1_12shot/}.
\item OLMo comparison: \file{results/olmo2_results.tar.gz} (runs \file{olmo2_7b_v3_1_0shot}, \file{olmo2_7b_v3_1_4shot}, \file{olmo2_7b_v3_1_12shot}, \file{olmo2_7b_v3_1_controls}); analysis and report: \file{results/olmo2_vs_pythia/}.
\end{itemize}

\begin{thebibliography}{9}
\bibitem{ren} Jie Ren et al. (2024). \emph{Identifying Semantic Induction Heads to Understand In-Context Learning}. Findings of ACL, 6916--6932. \url{https://aclanthology.org/2024.findings-acl.412/}
\bibitem{olsson} Catherine Olsson et al. (2022). \emph{In-context Learning and Induction Heads}. Transformer Circuits Thread. \url{https://transformer-circuits.pub/2022/in-context-learning-and-induction-heads/index.html}
\bibitem{framework} Nelson Elhage et al. (2021). \emph{A Mathematical Framework for Transformer Circuits}. Transformer Circuits Thread. \url{https://transformer-circuits.pub/2021/framework/index.html}
\bibitem{ioi} Kevin Wang et al. (2022). \emph{Interpretability in the Wild: a Circuit for Indirect Object Identification in GPT-2 Small}. \url{https://arxiv.org/abs/2211.00593}
\bibitem{patchguide} Stefan Heimersheim and Neel Nanda (2024). \emph{How to Use and Interpret Activation Patching}. \url{https://arxiv.org/abs/2404.15255}
\bibitem{olmo2} Team OLMo et al. (2024). \emph{2 OLMo 2 Furious}. \url{https://arxiv.org/abs/2501.00656}
\end{thebibliography}
\end{document}
